% Part: normal-modal-logic
% Chapter: syntax-and-semantics
% Section: relational-models

\documentclass[../../../include/open-logic-section]{subfiles}

\begin{document}

\olfileid[zh]{nml}{syn}{rel}

\olsection{关系模型}

正规模态逻辑语义的基本概念是\emph{关系模型}。它由一个世界集、这些世界之间的二元「可及关系」，以及一个决定哪些!!{propositional variable}在哪些世界处算作「真」的指派构成。

\begin{defn}
  基本模态语言的一个\emph{模型}是三元组 $\mModel{M}
  = \tuple{W, R, V}$，其中
  \begin{enumerate}
  \item $W$ 是一个「世界」的非空集，
  \item $R$ 是 $W$ 上的二元可及关系，并且
  \item $V$ 是一个给每个!!{propositional
    variable}~$p$ 指派一个可能世界集 $V(p)$ 的函数。
  \end{enumerate}
  当 $Rww'$ 成立时，我们说 $w'$ 从~$w$ \emph{可及}。当 $w \in
  V(p)$ 时，我们说 $p$ 在~$w$ 处\emph{为真}。
\end{defn}

关系语义的一大优点在于，模型可以用简单的图形表示，例如\olref{fig:simple} 中的模型。世界用节点表示；世界~$w'$ 从~$w$ 可及，当且仅当有一条从~$w$ 到~$w'$ 的箭头。此外，当 $w \in V(p)$ 时，我们用~$\mTrue{p}$ 标记相应的节点（世界），否则用~\mFalse{p} 标记。
\olref{fig:simple} 表示的模型具有 $W = \{w_1, w_2, w_3\}$，$R =
\{\tuple{w_1, w_2},\tuple{w_1, w_3}\}$，$V(p) = \{w_1, w_2\}$，以及
$V(q) = \{w_2\}$。

\begin{figure}
  \begin{center}
    \begin{tikzpicture}[modal]
      \node[world] (w1) [label={[align=right]right:\mTrue{p}\\ \mFalse{q}}]{$w_1$};
      \node[world] (w2) [label={[align=right]right:\mTrue{p}\\ \mTrue{q}},
        above right=of w1]{$w_2$};
      \node[world] (w3) [label={[align=right]right:\mFalse{p}\\ \mFalse{q}},
        below right=of w1] {$w_3$};
      \draw[->] (w1) to (w2);
      \draw[->] (w1) to (w3);
    \end{tikzpicture}
  \end{center}
  \caption{一个简单的模型。}
  \ollabel{fig:simple}
\end{figure}

\end{document}
